\section{The complete-kernel Hessian}\label{sec:hessian}

Let $P_\varepsilon=J_n+\varepsilon\delta$ with
$\delta\in\calT_n$.  By linearity,
\[
 K(P_\varepsilon)=K_0+\varepsilon\Delta.                    \tag{4.1}\label{4.1}
\]
Because every off-diagonal entry of $J_n$ is $1/(n-1)$, each fixed
$\delta$ admits $\varepsilon_0>0$ such that $P_\varepsilon$ remains strictly
positive off the diagonal for $|\varepsilon|<\varepsilon_0$.  Thus all
positive-kernel identities from Section~\ref{sec:duality} apply in the
neighborhood used below.  All entries are analytic there, and $K_0$ is
irreducible.  Put
\[
 q=H-c_0\one,\qquad
 G=(I-K_0+\one\nu_0)^{-1}.                                  \tag{4.2}\label{4.2}
\]

\subsection{Vanishing first variation}

Let $S$ average a function on $\mathcal Y_n$ over all active states of the
same rank $|B|$.  Permutation symmetry gives $SK_0=K_0S$.  A direct average
of either active move shows
\[
 S\Delta S=0.                                               \tag{4.3}\label{4.3}
\]
Indeed, a uniformly random $k$-subset of $V\setminus\{v\}$ contains a
sample from any row with mean probability $k/N$; after retargeting the
corresponding mean is $(k-1)/N$.  More explicitly, if $f_k$ is any rank
function, the two branchwise perturbations at $(B,v)$ are
\[
 \frac{f_{k+1}-f_k}{2}\sum_{i\notin B}\delta_{vi},\qquad
 \frac{f_{k-1}-f_k}{2k}
   \sum_{w\in B}\sum_{i\in B\setminus\{w\}}\delta_{wi}.    \tag{4.3a}\label{4.3a}
\]
Under the uniform rank-$k$ orbit their respective row sums are
\[
 \frac{N-k}{N}\sum_{i\ne v}\delta_{vi}=0,\qquad
 \frac{k-1}{N}\sum_{i\ne w}\delta_{wi}=0.                \tag{4.3b}\label{4.3b}
\]
This proves \eqref{4.3} without reversibility.  Because $q$ and $Gq$ are rank functions
and $\nu_0=\nu_0S$,
\[
 \nu_0\Delta Gq=\nu_0S\Delta SGq=0.                         \tag{4.4}\label{4.4}
\]

Stationary perturbation of the finite chain gives
\[
 \nu(P_\varepsilon)q
 =\varepsilon\nu_0\Delta Gq
  +\varepsilon^2\nu_0\Delta G\Delta Gq+O(\varepsilon^3).
                                                                    \tag{4.5}\label{4.5}
\]
Define
\[
 \Rtwo_{n}(\delta)=\nu_0\Delta G\Delta Gq.                    \tag{4.6}\label{4.6}
\]
Equations \eqref{3.12}, \eqref{3.17}, and \eqref{4.4}--\eqref{4.6}
already prove expansion \eqref{2.6}; it remains to prove that
$\Rtwo_{n}$ is positive definite on the full tangent space.

\subsection{The three irreducible tangent sectors}

For $\delta\in\calT_n$, let
\[
 s_j=\sum_i\delta_{ij},\qquad \sum_js_j=0,                   \tag{4.7}\label{4.7}
\]
and define the standard embedding
\[
 E(s)_{ij}=\frac{s_i+(n-1)s_j}{n(n-2)}\quad(i\ne j),
 \qquad E(s)_{ii}=0.                                       \tag{4.8}\label{4.8}
\]
Then $E(s)\one=0$ and its column-sum vector is $s$.  Consequently
$B=\delta-E(s)$ has zero row and column sums, and
\[
 \delta=E(s)+\frac{B+B^{\mathsf T}}2
                +\frac{B-B^{\mathsf T}}2.                    \tag{4.9}\label{4.9}
\]
The three terms are Frobenius-orthogonal.  Their dimensions are
\[
 n-1,\qquad \frac{n(n-3)}2,\qquad
 \frac{(n-1)(n-2)}2,                                      \tag{4.10}\label{4.10}
\]
which sum to $n(n-2)=\dim\calT_n$.  They are, respectively, the standard,
symmetric row/column-balanced, and antisymmetric balanced irreducible
$S_n$-modules.  At $n=3$ the symmetric module has dimension zero.

For completeness, the symmetric off-diagonal matrices form the two-subset
permutation module
$\mathbf1\oplus[n-1,1]\oplus[n-2,2]$; its row-balanced kernel is
$[n-2,2]$.  The antisymmetric balanced kernel is
$\bigwedge^2[n-1,1]=[n-2,1,1]$, while $E(s)$ realizes $[n-1,1]$.
Thus the three summands in \eqref{4.9} are pairwise nonisomorphic and occur
with multiplicity one, including the stated $n=3$ degeneration.

The quadratic form \eqref{4.6} is invariant under simultaneous relabeling of
targets and sources.  Multiplicity-freeness therefore makes the three spaces
orthogonal for $\Rtwo_{n}$ as well, and $\Rtwo_{n}$ acts by one scalar on each.
The balanced subspace is exactly
\[
 \calB_n=\calS_n\oplus\calA_n,                              \tag{4.10b}\label{4.10b}
\]
where $\calS_n$ and $\calA_n$ denote the symmetric- and
antisymmetric-balanced summands in \eqref{4.9}.  The standard module $E(s)$
is precisely the complementary column-imbalance direction.

\begin{theorem}[sector positivity]\label{thm:sectors}
The standard scalar is positive for every $n\ge3$, the
symmetric-balanced scalar is positive for every $n\ge4$, and the
antisymmetric-balanced scalar is positive for every $n\ge3$.
\end{theorem}

The proof is in Appendix~\ref{app:sectors}.  We summarize its architecture.
For the standard sector, a signed three-channel quotient separates positive
and negative orientations.  Schur elimination turns each completed negative
excursion into a nonnegative re-entry operator.  Exact first-phase barriers
and a uniform row-sum contraction control the entire alternating resolvent.
For the antisymmetric sector, a one-feature reduction sends a rank Poisson
gradient through a positive decreasing recurrence; the second perturbation
is then a sum of nonnegative sampling-without-replacement terms, at least one
strict.  For the symmetric sector, the exact quotient has a good/bad block form
\[
 \begin{pmatrix}S&C\\-D&Q\end{pmatrix}                     \tag{4.11}\label{4.11}
\]
with $S,C,D,Q$ nonnegative.  Schur elimination converts every completed bad
excursion into one factor of a nonnegative return operator $A$.  Explicit
positive phase barriers dominate the alternating resolvent $(I+A)^{-1}$.
The computer-assisted part is finite and exact: rational stationary solves
cover $3\le N\le39$, and rational phase-margin checks cover
$40\le N\le287$.  Coefficientwise polynomial and discriminant certificates
prove the remaining range $N\ge288$.  Appendix~\ref{app:sectors} states the
finite computation as a formal lemma and identifies its exact verifier.

For orientation, the first sector eigenvalues are
\[
\begin{array}{c|ccc}
n&\lambda_{\rm std}&\lambda_{\rm sym}&\lambda_{\rm anti}\\ \hline
3&1/11&---&1/9\\
4&87/640&3/208&57/640\\
5&8585/57314&359/26660&143/2100.
\end{array}                                                  \tag{4.12}\label{4.12}
\]
These are the Frobenius-normalized quadratic eigenvalues; they are checks,
not a finite extrapolation.

\begin{proof}[Proof of Theorem~\ref{thm:local}]
Theorem~\ref{thm:sectors} and the decomposition \eqref{4.9} imply
$\Rtwo_{n}(\delta)>0$ for every nonzero $\delta\in\calT_n$.  Equations
\eqref{3.12} and \eqref{4.5} give
\[
 \frac1{n\rho_{\dB}(P_\varepsilon,2)}
 =c_0+\varepsilon^2\Rtwo_{n}(\delta)+O(\varepsilon^3).         \tag{4.13}\label{4.13}
\]
Differentiation yields \eqref{2.8}.  To obtain a neighborhood rather than
only directional statements, restrict the analytic fixation map to a small
Frobenius ball around $J_n$ in the affine space $J_n+\calT_n$, chosen inside the
strictly positive off-diagonal region.  The smallest
$\Rtwo_{n}$-eigenvalue on the Frobenius unit sphere is positive.  Taylor's theorem with
a uniform third-derivative bound on a smaller compact ball makes the
negative quadratic term dominate the remainder.  Hence every other kernel
in that ball has strictly smaller fixation probability.
\end{proof}

\begin{remark}
The neighborhood in the last paragraph depends on $n$.  No statement here
prevents a population sequence from approaching a singular face, or from
remaining a nonperturbative distance from $J_n$ as $n\to\infty$.
\end{remark}
